
\problem{1}
У Вани есть 10 шорт и 10 футболок. Сможет ли мальчик все лето каждый день ходить гулять так, чтобы его наряд не повторился?

\textbf{Ответ:} да.

\textbf{Решение.}
Каждый день он может выбрать шорты одним из 10 способов. Для каждых шорт у него еще есть 10 способов выбрать футболку. Всего 100 комбинаций. 100 дней это даже больше, чем «все лето».

\problem{2}
а) Каким числом способов можно поставить 7 человек в колонну?

\noindent
б) Сколько существует трехзначных чисел, состоящих из различных цифр?

\textbf{Ответ:} а) $7 * 6 * 5 * 4 * 3 * 2 * 1=7!=5040$ : б) $9 * 9 * 8=648$

\textbf{Решение.}
а) Сначала выберем, кого поставим во главе колонны (первым). Любого из 7 человек. Потом выберем того, кто будет стоять вторым. У нас осталось 6 человек на выбор, не важно кого мы назначили первым. На третье место будем выбирать из 5 человек, на четвертое место из 4 человек, на пятое из 3, на шестое из 2, на последнее оставшегося одного человека. Всего $7 * 6 * 5 * 4 * 3 * 2 * 1$ вариантов.

\noindent
б) первую цифру можно выбрать любую от 1 до 9 (9 вариантов). После выбора первой цифры выберем вторую. Это можно сделать 9 способами, т.к. одну из цифр, выбранную для первой, использовать нельзя. Не важно, какая выбрана цифра на первое место, для выбора второй все равно остается 9 вариантов. Третью цифру можно выбрать 8 способами, т.к. она должна быть не равна первой и второй (и эти две недоступные цифры различны между собой). Значит, все три цифры можно выбрать 9*9*8 способами.


\problem{3}
Сколько существует семизначных чисел, в записи которых не встречается нулей, и которые делятся на 4 ?

\textbf{Oтвет:} $18^* 9^5=1062882$.

Peшение 1. 
Число делится на 4, если число образованное двумя его последними цифрами делится на 4. Окончаний чисел из двух цифр, делящихся на 4, - 25 штук (всего окончаний чисел из двух последних цифр 100 (00, 01, ..., 99), на 4 делятся каждые четвертые), но у нас нет нулей. Из-за отсутствия нулей выпадают окончания $00,04,08,20,40,60,80$, всего 7. Итак, окончаний чисел из двух цифр без нулей, делящихся на 4, - 18 (25-7). Предыдущие 5 цифр можно выбирать произвольно, каждую 9 способами, т.к. нет нулей. Это
$9 * 9 * 9 * 9 * 9=95$ вариантов. Для каждого из этих 95 вариантов выбора первых пяти цифр последние две мы можем выбрать 18 способами, т.е. всего вариантов выбрать все семь цифр будет $9^5$ *18.

Решение 2. Первые 5 чифр выберем произвольно, это можно сделать 95 способами. Если мы взяли предпоследнюю (шестую) чифру нечетную (еще 5 способов), то последняя может быть только 2 и 4, т.е. таких чисел $95 * 5 * 2$. Если же шестая чифра была четной (4 способа), то последняя, может быть только 4 или 8, таких чисел 95*4*2. Всего чисел $9^5 * 5 * 2+9^5 * 4 * 2=9^5 * 2 *(5+4)=9^6 * 2$.



\problem{4}
а) Каким числом способов из 10 различных горшков можно выбрать один для фикуса, один для крокуса, и один для кактуса?

\noindent
б) А каким числом способов из 10 горшков можно выбрать три, чтобы посадить в них лук?

\textbf{Ответ:} а) $10 * 9 * 8=720$; б) $10 * 9 * 8 /(3 * 2 * 1)=120$.

\textbf{Решение.} 
а) Горшок для фикуса можно взять любой из
10, потом горшок для крокуса возьмем любой из 9 оставшихся (не важно, в какой именно горшок посадили фикус, все равно останется 9 на выбор), и, наконец, кактус посадим в любой из оставшихся 8 горшков. Всего 10*9*8 вариантов.

\noindent
б) Итак, в пункте «а» мы нашли 720 способов посадить фикус, крокус и кактус. Пусть теперь и фикус, и крокус, и кактус заменили на лук. Вопрос: сколько способов из тех 720 станут одинаковыми (раньше эти способы были различные потому, что в одном горшке рос фикус, а в другом кактус, а теперь и там и там лук)?

Наблюдаем 3 горшка и 3 растения в этих горшках. Первое растение либо фикус, либо кактус, либо крокус. Второе - любое из двух оставшихся растений, последнее растение только одно и выбрать его можно одним способом. Т.е. у любого варианта рассадки трех растений в выбранные 3 горшка есть $3 * 2 * 1=6$ разных вариантов; все эти 6 вариантов для лука или других трех одинаковых растений станут одинаковыми.

Значит, чтобы найти количество вариантов рассадки лука, нужно количество вариантов рассадки фикуса-крокуса-кактуса разделить на 6.

\problem{5}
Михаил Александрович хочет выбрать из 13 семиклассников для участия в матбое капитана, заместителя капитана и остальных 4 членов команды. Сколькими способами он это может сделать?



\textbf{Oтвет:} $13 * 12 * 11 * 10 * 9 * 8 /(4!)=51480$.

\textbf{Решение.} 
Капитана команды можно выбрать 13 способами. После этого заместителя команды можно выбрать 12. Осталось 11 человек и из них надо набрать 4 рядовых членов команды. Будем рассуждать аналогично задаче 4б. Если бы у всех них были разные роли, то человека на первую роль можно было бы выбрать 11 способами, на вторую 10, на третью 9, на четвертую 8. Но при этом мы каждую набор из 4 членов команды посчитали столько раз сколькими способами можно разделить между ними эти 4 роли. А это $4 * 3 * 2 * 1=4!$ способов. Значит, набрать 4 рядовых членов команды из 11 оставшихся можно 11*10*9*8/(4!) способами, откуда всего способов $13 * 12 * 11 * 10 * 9 * 8 /(4!)$.

\problem{6}
Сколькими способами из 13 человек можно выбрать шесть для участия в матбое?

\textbf{Oтвет:} $13 * 12 * 11 * 10 * 9 * 8 /(6!)=1716$.

\textbf{Решение.}
Задача решается аналогично предыдущей задаче. Первого человека можно выбрать 13 способами, второго 12 и так далее - шестого 8 способами - всего $13 * 12 * 11 * 10 * 9 * 8$.

Но от того, что мы одних и тех же 6 человек берем в разном порядке команда не меняется. А порядков, в которых можно взять 6 человек - 6! (аналогично задаче 2а) Значит, всего способов будет 13*12*11*10*9*8/(6!).

\problem{7}
Забор состоит из 11 досок. Какие-то пять досок этого забора нужно покрасить в синий цвет. Сколькими способами это можно сделать?

\textbf{Oтвет:} $11 * 10 * 9 * 8 * 7 /(5!)=462$.

\textbf{Решение.} 
Аналогично предыдущей задаче. Первую доску можно покрасить 11 способами, вторую 10 и т.д. пятую - 7 способами. Но от того, в каком порядке мы красили доски, результат не меняется. Порядков, в которых можно покрасить 5 досок - 5 ! (аналогично задаче 2a). Значит, всего способов будет 11*10*9*8*7/(5!).

\problem{8}
Есть 5 белых и 6 черных шариков. Сколькими способами их можно расположить в ряд?
\textbf{Oтвет:} $11 * 10 * 9 * 8 * 7 /(5!)=462$.

\textbf{Решение.} 
Расположить в ряд 5 белых и 6 черных шариков, это все равно, что положить 11 черных шариков, а потом выбрать 5 их них и покрасить в белый цвет. Из прошлой задачи мы знаем, что это можно сделать $11 * 10 * 9 * 8 * 7 /(5!)$ способами.


\problem{9}
Математическая черепаха стоит в левой нижней клетке доски $7 \times 15$. За один ход она может передвинуться на 1 клетку вправо или на 1 клетку вверх. Сколькими способами черепаха может добраться в правый верхний угол?

\textbf{Oтвет:} $20 * 19 * 18 * 17 * 16 * 15 /(6!)=38760$.

\textbf{Решение.} 
Чтобы попасть из левого нижнего угла в правый верхний угол черепаха должна сделать 6 ходов вверх и 14 ходов вправо (или наоборот, в зависимости от того как стоит доска). Способов добраться из левого нижнего угла в правый верхний, столько же, сколько способов упорядочить эти 20 ходов (6 вверх и 14 вправо). Это то же самое, что определить на каких местах из 20 происходят 6 ходов вверх. Первое место можно выбрать 20 способами, второе 19 и т.д. шестое 15. От порядка, в котором мы выбирали способы, ничего не зависит, значит надо поделить еще на 6!.

\problem{10}
Сколькими способами можно расставить 5 ладей на доске $13 \times 13$ так, чтобы они не били друг друга?

\textbf{Oтвет:} $(13 * 13) *(12 * 12) *(11 * 11) *(10 * 10) *(9 * 9) /(5!)=198764280$.

\textbf{Решение.}
Будем выставлять ладьи по одной. Первую ладью можно поставить 13*13 способами, так как все клетки свободны. Так как ладья бьет одну вертикаль и одну горизонталь, то свободными остались 12 вертикалей и 12 горизонталей. Новую ладью можно поставить 12*12 способами. Теперь свободны уже 11 вертикалей и 11 горизонталей - третью ладью можно поставить 11*11 способами. И так далее. Теперь нам нужно сократить варианты, где одна и та же расстановка ладей получена выставлением ладей в разном порядке. Этих порядков всего $5!$ так как мы расставляли 5 ладей. Вот и получается, что способов всего (13*13)*(12*12)*(11*11) *(10*10)*(9*9)/(5!).